(1) 无黏不可压平面流体占据 \(y>0\)，在 \((X,h)\) 有逆时针环量 \(\Gamma\ne0\) 的点涡，此外有均匀 \(U\mathbf e_x\) 流。用 \(y<0\) 的镜像涡建立流函数，导出全场速度和壁面速度，并证明壁面无穿透。
(2) 忽略真实点涡的自身奇异速度，求其准确平移速度、离壁距离的演化及世界线 \(X(t)\)。说明 \(h\to\infty\) 和 \(\Gamma\to-\Gamma\) 的极限。
(3) 在以点涡速度平移的坐标系中建立定常流函数。应用正确的定常Bernoulli积分，求相对远场压强的壁面压力分布；解释不能直接在实验室坐标中使用无非定常项的Bernoulli式。
(4) 直接积分整条壁的压力差，证明净法向压力增量严格为零，同时计算壁面扰动切向速度的积分，解释两者并不矛盾。
